\begin{table}[htbp]
\centering
\caption{The class-absent test is paired, and is reported here as such. Each of the 27 acquisitions is scored twice, once against the full computed library and once with its own toxicological class deleted, so the two conditions are one sample under two treatments. The second column counts the acquisitions that move in the expected direction; the third is the exact two-sided sign test over the non-tied pairs. The two AUC columns are the same rank statistic computed on the acquisition and on the agent mean. The vial labeled HN2 is a condition of its own and is excluded, as elsewhere.}
\label{tab:si-abstpaired}
\small
\begin{tabular}{lrrrr}
\toprule
Score & Moves as expected & Sign test & AUC (acquisition) & AUC (agent) \\
\midrule
top-1 cosine & 26/26 & $<10^{-7}$ & 0.809 & 0.83 \\
margin & 19/26 & 0.029 & 0.758 & 0.78 \\
peak-to-sidelobe & 24/27 & 0.0000 & 0.904 & 0.91 \\
library-normalized $z$ & 26/26 & $<10^{-7}$ & 0.826 & 0.88 \\
unmixing residual & 26/26 & $<10^{-7}$ & 0.781 & 0.80 \\
band-match fraction & 13/24 & 0.839 & 0.669 & 0.66 \\
\bottomrule
\end{tabular}

\begin{flushleft}\footnotesize On the paired statistic 5 of the six scores separate the two conditions at the 5\,\% level. Three of them -- the top-1 cosine, the library-normalized $z$ and the unmixing residual -- return every pair in the expected direction, and that is a property of the design rather than evidence: deleting a class removes candidates, so the class-absent pool is a strict subset of the full pool, the top-1 cosine is a maximum over that pool, the $z$ is an affine function of it, and the residual is a least-squares residual on a subset of the same basis, so none of the three can move the other way. The paired claim therefore rests on the two scores that are free to: the peak-to-sidelobe ratio and the margin. The AUC is retained as an effect size because it is what the open-set literature quotes. The two AUC columns differ by at most 0.049, so the rank statistic is close to but not independent of the resampling unit, which is an empirical property of these data and not a theorem. None of this changes the conclusion of Table~\ref{tab:si-absten}: the scores order the two conditions, and no threshold separates them without error. Produced by \texttt{code/v3\_abstention\_paired.py}.\end{flushleft}
\end{table}
